\documentclass[10pt,a4paper]{amsart}
\usepackage{amsmath,amssymb,amsthm,graphicx}
\usepackage{verbatim}
\newenvironment{reference}{\comment}{\endcomment}
%\newenvironment{reference}{}{}
\newenvironment{slogan}{\comment}{\endcomment}
\newenvironment{history}{\comment}{\endcomment}

% For commutative diagrams we use Xy-pic
\usepackage[all]{xy}

% We use 2cell for 2-commutative diagrams.
\xyoption{2cell}
% Theorem environments.
%
\theoremstyle{plain}
\newtheorem{theorem}[subsection]{Theorem}
\newtheorem{proposition}[subsection]{Proposition}
\newtheorem{lemma}[subsection]{Lemma}

\theoremstyle{definition}
\newtheorem{definition}[subsection]{Definition}
\newtheorem{example}[subsection]{Example}
\newtheorem{exercise}[subsection]{Exercise}
\newtheorem{situation}[subsection]{Situation}

\theoremstyle{remark}
\newtheorem{remark}[subsection]{Remark}
\newtheorem{remarks}[subsection]{Remarks}

\numberwithin{equation}{subsection}

% Macros
%
\def\lim{\mathop{\mathrm{lim}}\nolimits}
\def\colim{\mathop{\mathrm{colim}}\nolimits}
\def\Spec{\mathop{\mathrm{Spec}}}
\def\Hom{\mathop{\mathrm{Hom}}\nolimits}
\def\Ext{\mathop{\mathrm{Ext}}\nolimits}
\def\SheafHom{\mathop{\mathcal{H}\!\mathit{om}}\nolimits}
\def\SheafExt{\mathop{\mathcal{E}\!\mathit{xt}}\nolimits}
\def\Sch{\mathit{Sch}}
\def\Mor{\mathop{\mathrm{Mor}}\nolimits}
\def\Ob{\mathop{\mathrm{Ob}}\nolimits}
\def\Sh{\mathop{\mathit{Sh}}\nolimits}
\def\NL{\mathop{N\!L}\nolimits}
\def\CH{\mathop{\mathrm{CH}}\nolimits}
\def\proetale{{pro\text{-}\acute{e}tale}}
\def\etale{{\acute{e}tale}}
\def\QCoh{\mathit{QCoh}}
\def\Ker{\mathop{\mathrm{Ker}}}
\def\Im{\mathop{\mathrm{Im}}}
\def\Coker{\mathop{\mathrm{Coker}}}
\def\Coim{\mathop{\mathrm{Coim}}}

% Boxtimes
%
\DeclareMathSymbol{\boxtimes}{\mathbin}{AMSa}{"02}

%
% Macros for moduli stacks/spaces
%
\def\QCohstack{\mathcal{QC}\!\mathit{oh}}
\def\Cohstack{\mathcal{C}\!\mathit{oh}}
\def\Spacesstack{\mathcal{S}\!\mathit{paces}}
\def\Quotfunctor{\mathrm{Quot}}
\def\Hilbfunctor{\mathrm{Hilb}}
\def\Curvesstack{\mathcal{C}\!\mathit{urves}}
\def\Polarizedstack{\mathcal{P}\!\mathit{olarized}}
\def\Complexesstack{\mathcal{C}\!\mathit{omplexes}}
% \Pic is the operator that assigns to X its picard group, usage \Pic(X)
% \Picardstack_{X/B} denotes the Picard stack of X over B
% \Picardfunctor_{X/B} denotes the Picard functor of X over B
\def\Pic{\mathop{\mathrm{Pic}}\nolimits}
\def\Picardstack{\mathcal{P}\!\mathit{ic}}
\def\Picardfunctor{\mathrm{Pic}}
\def\Deformationcategory{\mathcal{D}\!\mathit{ef}}
\usepackage{cleveref}
\crefname{originallemma}{original lemma}{original lemmas}
\Crefname{originallemma}{Original lemma}{Original lemmas}
\makeatletter
\newlabel{lemma-integral-over-invariants}{{03BJ}{source}{}{sourcepointer-1}{}}
\newlabel{lemma-integral-over-invariants@cref}{{[originallemma][0][]03BJ}{[1][0][]source}{}{}{}}
\makeatother
\crefname{originallemma}{original lemma}{original lemmas}
\Crefname{originallemma}{Original lemma}{Original lemmas}
\makeatletter
\newlabel{lemma-invariants-base-change}{{03BK}{source}{}{sourcepointer-2}{}}
\newlabel{lemma-invariants-base-change@cref}{{[originallemma][0][]03BK}{[1][0][]source}{}{}{}}
\makeatother
\crefname{originallemma}{original lemma}{original lemmas}
\Crefname{originallemma}{Original lemma}{Original lemmas}
\makeatletter
\newlabel{lemma-points}{{03BL}{source}{}{sourcepointer-3}{}}
\newlabel{lemma-points@cref}{{[originallemma][0][]03BL}{[1][0][]source}{}{}{}}
\makeatother
\crefname{originalreference}{original reference}{original references}
\Crefname{originalreference}{Original reference}{Original references}
\makeatletter
\newlabel{limits-lemma-descend-opens}{{01Z4}{source}{}{sourcepointer-4}{}}
\newlabel{limits-lemma-descend-opens@cref}{{[originalreference][0][]01Z4}{[1][0][]source}{}{}{}}
\makeatother
\crefname{originalreference}{original reference}{original references}
\Crefname{originalreference}{Original reference}{Original references}
\makeatletter
\newlabel{limits-lemma-descend-isomorphism}{{081E}{source}{}{sourcepointer-5}{}}
\newlabel{limits-lemma-descend-isomorphism@cref}{{[originalreference][0][]081E}{[1][0][]source}{}{}{}}
\makeatother
\crefname{originalreference}{original reference}{original references}
\Crefname{originalreference}{Original reference}{Original references}
\makeatletter
\newlabel{morphisms-lemma-set-points-where-fibres-etale}{{0476}{source}{}{sourcepointer-6}{}}
\newlabel{morphisms-lemma-set-points-where-fibres-etale@cref}{{[originalreference][0][]0476}{[1][0][]source}{}{}{}}
\makeatother
\crefname{originalreference}{original reference}{original references}
\Crefname{originalreference}{Original reference}{Original references}
\makeatletter
\newlabel{algebra-lemma-finite-residue-extension-closed}{{00GA}{source}{}{sourcepointer-7}{}}
\newlabel{algebra-lemma-finite-residue-extension-closed@cref}{{[originalreference][0][]00GA}{[1][0][]source}{}{}{}}
\makeatother
\crefname{originalreference}{original reference}{original references}
\Crefname{originalreference}{Original reference}{Original references}
\makeatletter
\newlabel{more-algebra-lemma-integral-over-henselian-pair}{{09XK}{source}{}{sourcepointer-8}{}}
\newlabel{more-algebra-lemma-integral-over-henselian-pair@cref}{{[originalreference][0][]09XK}{[1][0][]source}{}{}{}}
\makeatother
\crefname{originalreference}{original reference}{original references}
\Crefname{originalreference}{Original reference}{Original references}
\makeatletter
\newlabel{more-algebra-lemma-change-ideal-henselian-pair}{{09XJ}{source}{}{sourcepointer-9}{}}
\newlabel{more-algebra-lemma-change-ideal-henselian-pair@cref}{{[originalreference][0][]09XJ}{[1][0][]source}{}{}{}}
\makeatother
\crefname{originalreference}{original reference}{original references}
\Crefname{originalreference}{Original reference}{Original references}
\makeatletter
\newlabel{more-algebra-lemma-characterize-henselian-pair}{{09XI}{source}{}{sourcepointer-10}{}}
\newlabel{more-algebra-lemma-characterize-henselian-pair@cref}{{[originalreference][0][]09XI}{[1][0][]source}{}{}{}}
\makeatother
\crefname{originalreference}{original reference}{original references}
\Crefname{originalreference}{Original reference}{Original references}
\makeatletter
\newlabel{morphisms-lemma-etale-permanence}{{02GW}{source}{}{sourcepointer-11}{}}
\newlabel{morphisms-lemma-etale-permanence@cref}{{[originalreference][0][]02GW}{[1][0][]source}{}{}{}}
\makeatother
\crefname{originalreference}{original reference}{original references}
\Crefname{originalreference}{Original reference}{Original references}
\makeatletter
\newlabel{descent-lemma-universally-injective-etale-open-immersion}{{02LC}{source}{}{sourcepointer-12}{}}
\newlabel{descent-lemma-universally-injective-etale-open-immersion@cref}{{[originalreference][0][]02LC}{[1][0][]source}{}{}{}}
\makeatother
\title{Etale base change of a finite locally free groupoid: invariants etale and U' = M' x\_M U}
\author{The Stacks Project Authors}
\date{stacks-project:groupoids.tex:local-snapshot-20261012; GNU FDL 1.2 or later}
\begin{document}
\maketitle
\noindent\textit{Editorial note.} All delivered mathematical text is the author's own, copied verbatim from the pinned Stacks Project chapter file groupoids.tex; no statement or proof step was rewritten, reordered or removed.\par\medskip
\noindent\textit{Editorial note.} Span recomputation by label pairing (LANE\_BRIEF 10.2): in this chapter the author's proof environment follows the statement environment as a standalone block; the spec covers the complete environment pair (statement span, proof span) as recorded in the pool candidate.\par\medskip
\noindent\textit{Editorial note.} Occupation scan: this chapter's pool has 47 of 128 labels already delivered by the S7 line (S7-gp-* registrations, statement\_anchor-prose scan; withdrawn cards excluded); none of the 11 accepted labels of this batch is among them (gp-occupied.json, 2026-10-14).\par\medskip
\ifdefined\hypertarget\hypertarget{sourcepointer-1}{}\fi
\noindent\textit{Source reference.} Original lemma 03BJ: sources/stacks-chow/tags.txt, source lines 4350--4350. Cross-chapter standard prerequisite cited by the author inside the counted span; it is used as a pointer and its own proof is not part of this item.\par\medskip
\ifdefined\hypertarget\hypertarget{sourcepointer-2}{}\fi
\noindent\textit{Source reference.} Original lemma 03BK: sources/stacks-chow/tags.txt, source lines 4351--4351. Cross-chapter standard prerequisite cited by the author inside the counted span; it is used as a pointer and its own proof is not part of this item.\par\medskip
\ifdefined\hypertarget\hypertarget{sourcepointer-3}{}\fi
\noindent\textit{Source reference.} Original lemma 03BL: sources/stacks-chow/tags.txt, source lines 4352--4352. Cross-chapter standard prerequisite cited by the author inside the counted span; it is used as a pointer and its own proof is not part of this item.\par\medskip
\ifdefined\hypertarget\hypertarget{sourcepointer-4}{}\fi
\noindent\textit{Source reference.} Original reference 01Z4: sources/stacks-chow/tags.txt, source lines 2579--2579. Cross-chapter standard prerequisite cited by the author inside the counted span; it is used as a pointer and its own proof is not part of this item.\par\medskip
\ifdefined\hypertarget\hypertarget{sourcepointer-5}{}\fi
\noindent\textit{Source reference.} Original reference 081E: sources/stacks-chow/tags.txt, source lines 10214--10214. Cross-chapter standard prerequisite cited by the author inside the counted span; it is used as a pointer and its own proof is not part of this item.\par\medskip
\ifdefined\hypertarget\hypertarget{sourcepointer-6}{}\fi
\noindent\textit{Source reference.} Original reference 0476: sources/stacks-chow/tags.txt, source lines 5497--5497. Cross-chapter standard prerequisite cited by the author inside the counted span; it is used as a pointer and its own proof is not part of this item.\par\medskip
\ifdefined\hypertarget\hypertarget{sourcepointer-7}{}\fi
\noindent\textit{Source reference.} Original reference 00GA: sources/stacks-chow/tags.txt, source lines 600--600. Cross-chapter standard prerequisite cited by the author inside the counted span; it is used as a pointer and its own proof is not part of this item.\par\medskip
\ifdefined\hypertarget\hypertarget{sourcepointer-8}{}\fi
\noindent\textit{Source reference.} Original reference 09XK: sources/stacks-chow/tags.txt, source lines 12533--12533. Cross-chapter standard prerequisite cited by the author inside the counted span; it is used as a pointer and its own proof is not part of this item.\par\medskip
\ifdefined\hypertarget\hypertarget{sourcepointer-9}{}\fi
\noindent\textit{Source reference.} Original reference 09XJ: sources/stacks-chow/tags.txt, source lines 12532--12532. Cross-chapter standard prerequisite cited by the author inside the counted span; it is used as a pointer and its own proof is not part of this item.\par\medskip
\ifdefined\hypertarget\hypertarget{sourcepointer-10}{}\fi
\noindent\textit{Source reference.} Original reference 09XI: sources/stacks-chow/tags.txt, source lines 12531--12531. Cross-chapter standard prerequisite cited by the author inside the counted span; it is used as a pointer and its own proof is not part of this item.\par\medskip
\ifdefined\hypertarget\hypertarget{sourcepointer-11}{}\fi
\noindent\textit{Source reference.} Original reference 02GW: sources/stacks-chow/tags.txt, source lines 3243--3243. Cross-chapter standard prerequisite cited by the author inside the counted span; it is used as a pointer and its own proof is not part of this item.\par\medskip
\ifdefined\hypertarget\hypertarget{sourcepointer-12}{}\fi
\noindent\textit{Source reference.} Original reference 02LC: sources/stacks-chow/tags.txt, source lines 3403--3403. Cross-chapter standard prerequisite cited by the author inside the counted span; it is used as a pointer and its own proof is not part of this item.\par\medskip
\begin{lemma}
\label{lemma-etale}
Let $S$ be a scheme. Let $f : (U', R', s', t') \to (U, R, s, t, c)$ be a
morphism of groupoid schemes over $S$.
\begin{enumerate}
\item $U$, $R$, $U'$, $R'$ are affine,
\item $s, t, s', t'$ are finite locally free,
\item the diagrams
$$
\xymatrix{
R' \ar[d]_{s'} \ar[r]_f & R \ar[d]^s \\
U' \ar[r]^f & U
}
\quad
\quad
\xymatrix{
R' \ar[d]_{t'} \ar[r]_f & R \ar[d]^t \\
U' \ar[r]^f & U
}
\quad
\quad
\xymatrix{
G' \ar[d] \ar[r]_f & G \ar[d] \\
U' \ar[r]^f & U
}
$$
are cartesian where $G$ and $G'$ are the stabilizer group schemes, and
\item $f : U' \to U$ is \'etale.
\end{enumerate}
Then the map $C \to C'$ from the $R$-invariant functions on $U$
to the $R'$-invariant functions on $U'$ is \'etale and
$U' = \Spec(C') \times_{\Spec(C)} U$.
\end{lemma}

\begin{proof}
Set $M = \Spec(C)$ and $M' = \Spec(C')$.
Write $U = \Spec(A)$, $U' = \Spec(A')$, $R = \Spec(B)$, and
$R' = \Spec(B')$. We will use the results of
Lemmas \ref{lemma-integral-over-invariants},
\ref{lemma-invariants-base-change}, and
\ref{lemma-points}
without further mention.

\medskip\noindent
Assume $C$ is a strictly henselian local ring. Let $p \in M$
be the closed point and let $p' \in M'$ map to $p$.
Claim: in this case there is a disjoint union decomposition
$(U', R', s', t', c') = (U, R, s, t, c) \amalg (U'', R'', s'', t'', c'')$
over $(U, R, s, t, c)$ such that for the corresponding
disjoint union decomposition $M' = M \amalg M''$ over $M$
the point $p'$ corresponds to $p \in M$.

\medskip\noindent
The claim implies the lemma. Suppose that $M_1 \to M$ is a flat morphism
of affine schemes. Then we can base change everything to $M_1$
without affecting the hypotheses (1) -- (4).
From Lemma \ref{lemma-invariants-base-change}
we see $M_1$, resp.\ $M_1'$ is the spectrum of the
$R_1$-invariant functions on $U_1$,
resp.\ the $R'_1$-invariant functions on $U'_1$.
Suppose that $p' \in M'$ maps to $p \in M$.
Let $M_1$ be the spectrum of the strict henselization of
$\mathcal{O}_{M, p}$ with closed point $p_1 \in M_1$.
Choose a point $p'_1 \in M'_1$ mapping to $p_1$ and $p'$.
From the claim we get
$$
(U'_1, R'_1, s'_1, t'_1, c'_1) =
(U_1, R_1, s_1, t_1, c_1) \amalg
(U''_1, R''_1, s''_1, t''_1, c''_1)
$$
and correspondingly $M'_1 = M_1 \amalg M''_1$ as a scheme over $M_1$.
Write $M_1 = \Spec(C_1)$ and write $C_1 = \colim C_i$ as a filtered
colimit of \'etale $C$-algebras. Set $M_i = \Spec(C_i)$.
The $M_1 = \lim M_i$ and similarly for the other schemes.
By Limits, Lemmas \ref{limits-lemma-descend-opens} and
\ref{limits-lemma-descend-isomorphism}
we can find an $i$ such that
$$
(U'_i, R'_i, s'_i, t'_i, c'_i) =
(U_i, R_i, s_i, t_i, c_i) \amalg
(U''_i, R''_i, s''_i, t''_i, c''_i)
$$
We conclude that $M'_i = M_i \amalg M''_i$. In particular
$M' \to M$ becomes \'etale at a point over $p'$ after an
\'etale base change. This implies that $M' \to M$ is \'etale at $p'$
(for example by Morphisms, Lemma
\ref{morphisms-lemma-set-points-where-fibres-etale}).
We will prove $U' \cong M' \times_M U$ after we prove the claim.

\medskip\noindent
Proof of the claim. Observe that $U_p$ and $U'_{p'}$ have finitely many points.
For $u \in U_p$ we have $\kappa(u)/\kappa(p)$ is algebraic,
hence $\kappa(u)$ is separably closed.
As $U' \to U$ is \'etale, we conclude the morphism $U'_{p'} \to U_p$
induces isomorphisms on residue field extensions.
Let $u' \in U'_{p'}$ with image $u \in U_p$.
By assumption (3) the morphism of scheme theoretic fibres
$(s')^{-1}(u') \to s^{-1}(u)$,
$(t')^{-1}(u') \to t^{-1}(u)$, and
$G'_{u'} \to G_u$ are isomorphisms. Observing that $U_p = t(s^{-1}(u))$
(set theoretically) we conclude that the points of $U'_{p'}$
surject onto the points of $U_p$.
Suppose that $u'_1$ and $u'_2$ are points of $U'_{p'}$ mapping
to the same point $u$ of $U_p$. Then there exists a point
$r' \in R'_{p'}$ with $s'(r') = u'_1$ and $t'(r') = u'_2$.
Consider the two towers of fields
$$
\kappa(r')/\kappa(u'_1)/\kappa(u)/\kappa(p) \quad
\kappa(r')/\kappa(u'_2)/\kappa(u)/\kappa(p)
$$
whose ``ends'' are the same as the two ``ends'' of the two towers
$$
\kappa(r')/\kappa(u'_1)/\kappa(p')/\kappa(p) \quad
\kappa(r')/\kappa(u'_2)/\kappa(p')/\kappa(p)
$$
These two induce the same maps $\kappa(p') \to \kappa(r')$ as
$(U'_{p'}, R'_{p'}, s', t', c')$ is a groupoid over $p'$.
Since $\kappa(u)/\kappa(p)$ is purely inseparable,
we conclude that the two induced maps
$\kappa(u) \to \kappa(r')$ are the same.
Therefore $r'$ maps to a point of the fibre $G_u$.
By assumption (3) we conclude that $r' \in (G')_{u'_1}$.
Namely, we may think of $G$ as a closed subscheme of $R$
viewed as a scheme over $U$ via $s$ and use that
the base change to $U'$ gives $G' \subset R'$.
In particular we have $u'_1 = u'_2$.
We conclude that $U'_{p'} \to U_p$ is a bijective
map on points inducing isomorphisms on residue fields.
It follows that $U'_{p'}$ is a finite set of closed points
(Algebra, Lemma \ref{algebra-lemma-finite-residue-extension-closed})
and hence $U'_{p'}$ is closed in $U'$.
Let $J' \subset A'$ be the radical ideal cutting out $U'_{p'}$
set theoretically.

\medskip\noindent
Second part proof of the claim.
Let $\mathfrak m \subset C$ be the maximal ideal.
Observe that $(A, \mathfrak m A)$ is a henselian pair by
More on Algebra, Lemma \ref{more-algebra-lemma-integral-over-henselian-pair}.
Let $J = \sqrt{\mathfrak m A}$.
Then $(A, J)$ is a henselian pair
(More on Algebra, Lemma \ref{more-algebra-lemma-change-ideal-henselian-pair})
and the \'etale ring map
$A \to A'$ induces an isomorphism $A/J \to A'/J'$
by our deliberations above.
We conclude that $A' = A \times A''$ by
More on Algebra, Lemma \ref{more-algebra-lemma-characterize-henselian-pair}.
Consider the corresponding disjoint union
decomposition $U' = U \amalg U''$. The open $(s')^{-1}(U)$ is the
set of points of $R'$ specializing to a point of $R'_{p'}$.
Similarly for $(t')^{-1}(U)$. Similarly we have
$(s')^{-1}(U'') = (t')^{-1}(U'')$ as this is the set of
points which do not specialize to $R'_{p'}$.
Hence we obtain a disjoint union decomposition
$$
(U', R', s', t', c') =
(U, R, s, t, c) \amalg
(U'', R'', s'', t'', c'')
$$
This immediately gives $M' = M \amalg M''$ and the proof of the claim
is complete.

\medskip\noindent
We still have to prove that the canonical map $U' \to M' \times_M U$
is an isomorphism. It is an \'etale morphism
(Morphisms, Lemma \ref{morphisms-lemma-etale-permanence}).
On the other hand, by base changing to strictly henselian local rings
(as in the third paragraph of the proof) and using the bijectivity
$U'_{p'} \to U_p$ established in the course of the proof of the claim,
we see that $U' \to M' \times_M U$ is universally bijective
(some details omitted). However, a universally bijective
\'etale morphism is an isomorphism
(Descent, Lemma \ref{descent-lemma-universally-injective-etale-open-immersion})
and the proof is complete.
\end{proof}

\end{document}
